\section{Business, Business Models and Their Characteristics}
% ─────────────────────────────────────────────────────────────

Businesses rely on transactions. Simply, a transaction is an economic event that
involves the exchange of either goods, services, or cash between two or more
parties.\footcite{shopify_transaction} For a transaction to occur:

\begin{enumerate}
    \item The buyer and seller must agree on the terms of the exchange.
    \item The buyer must offer something of value to the seller.
\end{enumerate}

A \textbf{business model} is a ``high-level plan for profitably operating a business
in a specific marketplace.''

% ─────────────────────────────────────────────────────────────
\subsection{Components of a Business Model}
% ─────────────────────────────────────────────────────────────

\begin{itemize}
    \item Value Proposition.
    \item Target Audience.
    \item Products or Services.
    \item Anticipated Expenses.
    \item Primary levels: Pricing and Costs. Their difference.
\end{itemize}

% ─────────────────────────────────────────────────────────────
\subsection{Types of Business Models}
% ─────────────────────────────────────────────────────────────

\begin{itemize}
    \item Retailers.
    \item Manufacturers.
    \item Fee-for-Service.
    \item Subscription.
    \item Freemium.
    \item Bundling.
    \item Marketplace.
    \item Affiliate.
    \item Razor Blade.
    \item Reverse Razor Blade.
    \item Franchise.
    \item Pay-as-You-Go.
    \item Brokerage.
\end{itemize}

% ─────────────────────────────────────────────────────────────
\subsection{Evaluating a Business Model}
% ─────────────────────────────────────────────────────────────

\subsubsection{Gross Profit}

% ── Business Model Canvas ────────────────────────────────────
\subsubsection{Business Model Canvas (9 Blocks)}

The Business Model Canvas is organized across three dimensions:

% ── Desirability ─────────────────────────────────────────────
\paragraph{Desirability (Do customers want it?)}\mbox{}\\[4pt]

\textbf{Customer Segments}
\begin{itemize}
    \item The heart of any model.
    \item Grouped based on needs, behaviors, and attributes.
    \item Who are we creating value for? Who are the most important customers?
\end{itemize}

\medskip

\textbf{Value Propositions}
\begin{itemize}
    \item The unique bundle of products/services that can solve customer problems.
    \item Value can be:
    \begin{itemize}
        \item \textit{Quantitative:} Price, Speed of Service.
        \item \textit{Qualitative:} Design, Customer Experience.
    \end{itemize}
    \item Why should customers choose us over competitors?
\end{itemize}

\medskip

\textbf{Channels}
\begin{itemize}
    \item How a company communicates with and reaches customer segments.
    \item Through which channels do customers want to be reached? How are we reaching
    them now?
\end{itemize}

\medskip

\textbf{Customer Relationships}
\begin{itemize}
    \item Relationship established with customer segment.
    \item Can range from personal to fully automated.
    \item Serves three core goals:
    \begin{itemize}
        \item Customer acquisition.
        \item Customer retention.
        \item Up-selling.
    \end{itemize}
\end{itemize}

% ── Visibility ───────────────────────────────────────────────
\paragraph{Visibility (Was it worth it?)}\mbox{}\\[4pt]

% ── Feasibility ──────────────────────────────────────────────
\paragraph{Feasibility (Can we deliver it?)}\mbox{}\\[4pt]

% ── Adaptability ─────────────────────────────────────────────
\paragraph{Adaptability (Can the idea survive and evolve?)}\mbox{}\\[4pt]

% ─────────────────────────────────────────────────────────────
\subsection{Business Model Analysis Template}
% ─────────────────────────────────────────────────────────────

\begin{itemize}
    \item Who are the parties involved?
    \item What value is typically exchanged, and how?
    \item What enables the model to work?
    \item What are limitations or risks?
    \item How does it differ from similar models?
\end{itemize}

% ─────────────────────────────────────────────────────────────
\subsection{Business Model Types}
% ─────────────────────────────────────────────────────────────

% ── C2C ──────────────────────────────────────────────────────
\subsubsection{Consumer-to-Consumer (C2C)}

A person sells a personal item or service to another person, usually through eBay,
Kijiji, or Facebook Marketplace.

\medskip

The facilitating companies such as eBay, Kijiji, or Facebook Marketplace provide
infrastructure, trust, and payment processing.

% ── B2B ──────────────────────────────────────────────────────
\subsubsection{Business-to-Business (B2B)}

A business sells a product or service that helps another business operate, produce, or
grow.

\medskip

Driven by ROI (return on investment). Why?

\medskip

This process is highly difficult and tedious. Why?

\medskip

Requires deep ICPC expertise. Without strong knowledge, B2B can be a waste of time.

\medskip

The call structure should generally follow this outline:
\begin{enumerate}
    \item Pre-call Prep.
    \item First 3 minutes.
    \item Discovery.
    \item Present fit.
    \item Handle Objections.
    \item Gain commitment.
\end{enumerate}

% ── B2C ──────────────────────────────────────────────────────
\subsubsection{Business-to-Consumer (B2C)}

% ── D2C ──────────────────────────────────────────────────────
\subsubsection{Direct-to-Consumer (D2C)}

% ── SaaS ─────────────────────────────────────────────────────
\subsubsection{Software-as-a-Service (SaaS)}

\begin{enumerate}
    \item Delivering centrally hosted applications over the internet on demand.
    \item It is a system where a business pays a subscription for digital access to
    software. Assertive is a great example of this.
    \item Some top SaaS companies include Apple, Microsoft, Google (owned by Alphabet),
    SAP, IBM, etc.
    \item Customers' lifetime is annuity. It is the yearly renewal of their account.
    \item Companies typically send out RFPs (request for protocols) when they need a
    service. They send out the details to multiple vendors and let them compete by
    submitting bids. It is a highly competitive cycle.
    \item The end goal for a tech startup is either to go public (IPO) or to be
    acquired or join a merger.
    \item Software in the 1980s was sold as a standalone product that you purchased and
    owned outright; came in the form of disks, along with licenses that gave the buyer
    access to some degree.
\end{enumerate}

% ── B2B SaaS ─────────────────────────────────────────────────
\subsubsection{Business-to-Business SaaS (B2B SaaS)}

% ── B2C SaaS ─────────────────────────────────────────────────
\subsubsection{Business-to-Consumer SaaS (B2C SaaS)}

% ── IaaS ─────────────────────────────────────────────────────
\subsubsection{Infrastructure-as-a-Service (IaaS)}

% ── PaaS ─────────────────────────────────────────────────────
\subsubsection{Platform-as-a-Service (PaaS)}

% ── B2G ──────────────────────────────────────────────────────
\subsubsection{Business-to-Government (B2G)}